% This LaTeX was auto-generated from MATLAB code.
% To make changes, update the MATLAB code and export to LaTeX again.

\documentclass{article}

\usepackage[utf8]{inputenc}
\usepackage[T1]{fontenc}
\usepackage{lmodern}
\usepackage{graphicx}
\usepackage{color}
\usepackage{hyperref}
\usepackage{amsmath}
\usepackage{amsfonts}
\usepackage{epstopdf}
\usepackage[table]{xcolor}
\usepackage{matlab}

\sloppy
\epstopdfsetup{outdir=./}
\graphicspath{ {./pbb1_images/} }

\begin{document}

\begin{par}
\begin{flushleft}
Examen Calcul numeric - Fabian-Ionuț Martin, grupa 234/2
\end{flushleft}
\end{par}

\begin{matlabcode}
syms x
f = @(x) (x*exp(x)-1);
f_sym = f(x)
\end{matlabcode}
\begin{matlabsymbolicoutput}
f\_sym = 

\hskip1em $\displaystyle x\,{\mathrm{e}}^x -1$
\end{matlabsymbolicoutput}
\begin{matlabcode}
fd = derivataDeOrdinNPlusUnu(f,0)
\end{matlabcode}
\begin{matlabsymbolicoutput}
fd = 

\hskip1em $\displaystyle {\mathrm{e}}^x +x\,{\mathrm{e}}^x $
\end{matlabsymbolicoutput}
\begin{matlabcode}
fd = matlabFunction(fd);
fdd = derivataDeOrdinNPlusUnu(f,1)
\end{matlabcode}
\begin{matlabsymbolicoutput}
fdd = 

\hskip1em $\displaystyle 2\,{\mathrm{e}}^x +x\,{\mathrm{e}}^x $
\end{matlabsymbolicoutput}
\begin{matlabcode}
fdd = matlabFunction(fdd);

eroare = 1e-10;
maxIteratii = 200;
[z, ni] = Aproximare(f,fd,fdd,0,eroare, maxIteratii);


fprintf("Metoda lui Newton:");
\end{matlabcode}
\begin{matlaboutput}
Metoda lui Newton:
\end{matlaboutput}
\begin{matlabcode}
fprintf("f(x) = x*exp(x)-1");
\end{matlabcode}
\begin{matlaboutput}
f(x) = x*exp(x)-1
\end{matlaboutput}
\begin{matlabcode}
fprintf("Numărul de iterații efectuate: %d", ni);
\end{matlabcode}
\begin{matlaboutput}
Numărul de iterații efectuate: 4
\end{matlaboutput}
\begin{matlabcode}
fprintf("Valoarea este = %.10f", z);
\end{matlabcode}
\begin{matlaboutput}
Valoarea este = 0.5671432904
\end{matlaboutput}
\begin{matlabcode}
fprintf("Valoarea aproximata de Matlab prin fzero: %.10f", fzero(f, 0));
\end{matlabcode}
\begin{matlaboutput}
Valoarea aproximata de Matlab prin fzero: 0.5671432904
\end{matlaboutput}
\begin{matlabcode}
fprintf("Valoarea aproximata de Matlab prin solve: %.10f", solve(f_sym));
\end{matlabcode}
\begin{matlaboutput}
Valoarea aproximata de Matlab prin solve: 0.5671432904
\end{matlaboutput}


\begin{par}
\begin{flushleft}
Proceduri folosite:
\end{flushleft}
\end{par}

\begin{par}
\begin{flushleft}
1) pentru aproximare
\end{flushleft}
\end{par}

\begin{matlabcode}
function [z,ni]=Aproximare(f,fd,fdd,x0,eroarea,maxIteratii)
    if nargin < 6, maxIteratii=50; end
    xp=x0(:);   %x precedent
    for k=1:maxIteratii
        xc=xp-f(xp)/sqrt(fd(xp)^2-f(xp)*fdd(xp)); 
        if abs(xc-xp)<eroarea
            z=xc; %succes
            ni=k;
            return
        end
        xp=xc;
    end
    error('S-a depasit numarul maxim de iteratii');
end
\end{matlabcode}

\begin{par}
\begin{flushleft}
2) pentru derivate
\end{flushleft}
\end{par}

\begin{matlabcode}
function [dfN] = derivataDeOrdinNPlusUnu(f, n) 
    syms x; 
    % Calculează derivata funcției f în raport cu x
    for i=1:n+1
        dfN = diff(f, x);
        f = dfN;
    end
end
\end{matlabcode}

\end{document}
